\section{Extended Parameters and Benchmarks}
\label{app:ext-params}

This appendix provides extended parameter discussion and additional benchmark data.

\subsection{Admissibility checklist}
A parameter tuple is \emph{admissible} for ARC-SPRUCE if it simultaneously satisfies:
\begin{enumerate}[leftmargin=1.25em]
  \item \textbf{Bound vs.\ modulus:} \(\BoundMsg \ll q\) so membership checks on \([ -\BoundMsg,\BoundMsg]\) do not wrap modulo \(q\).
  \item \textbf{Commitment hiding:} the commitment-only randomness \(\bar r\) satisfies the leftover-hash condition of Lemma~\ref{lem:ajtai-hiding}.
  \item \textbf{Target hiding:} the long BB-tran randomizer \(r_0\in\Rq^{\ell_{r_0}}\) satisfies the cyclotomic-LHL condition of Lemma~\ref{lem:htran-lhl}.
  \item \textbf{Sampler hiding and shortness:} Gaussian widths exceed the required smoothing threshold and the signature
  bound \(\BoundSig\) holds with overwhelming probability.
  \item \textbf{Rational-hash guard:} denominator invertibility succeeds with negligible failure probability under the
  holder randomness distribution.
  \item \textbf{SmallWood degree/soundness:} the DECS/LVCS degree bounds and PIOP soundness parameters meet the target
  security \(\lambda\).
  \item \textbf{PRF entropy:} truncation length \(\ell_{\prftag}\) yields at least \(\lambda\) bits of tag entropy.
\end{enumerate}

\subsection{Smoothing and sampler budgets}
Let \(\eta_\varepsilon(\Lambda)\) denote the smoothing parameter of a lattice \(\Lambda\). A sufficient condition for
statistical hiding in Gaussian preimage sampling is to choose Gaussian widths above a floor
\(\sigma_{\min}\ge \eta_\varepsilon(\Lambda_\perp(A))\), with slack (Appendix~\ref{app:gaussian-sampler}).
For the ambient lattice \(\Z^{2N}\) one may use the classical bound
\[
  \eta_\varepsilon(\Z^{2N})
  \,=\,
  \sqrt{\ln\!\bigl(4N(1+2^\lambda)\bigr)/\pi}.
\]
In our tuning, trapdoor quality \(\alpha\) and acceptance parameters are chosen so that the output shortness bound
\(\BoundSig\) holds with overwhelming probability.

\subsection{Example lattice parameters}
Table~\ref{tab:lattice-params} records one concrete instantiation compatible with the sizing discussion in
Section~\ref{sec:parameters}.

\begin{table}[H]
\centering
\renewcommand{\arraystretch}{1.15}
\begin{tabular}{@{}lll@{}}
\toprule
\textbf{Knob / Result} & \textbf{Symbol} & \textbf{Value (example)} \\
\midrule
Ring degree & \(N\) & \(1024\) \\
Modulus & \(q\) & \(1038337\) \\
Security target & \(\lambda\) & \(128\) \\
\midrule
Trapdoor quality & \(\alpha\) & \(\approx 1.23\) (measured) \\
Gaussian width scale & \(\sigma\) & tuned per-slot (Appendix~\ref{app:gaussian-sampler}) \\
Signature bound & \(\BoundSig\) & derived from sampler tail bound \\
\bottomrule
\end{tabular}
\caption{Example lattice parameters (illustrative; see Appendix~\ref{app:gaussian-sampler} for the sampling interface and
tuning).}
\label{tab:lattice-params}
\end{table}

\subsection{LHL design points for the long $r_0$-randomizer}
For the measured parameter scale \(N=1024\), \(q=1038337\), and \(\lambda=128\), the one-coordinate target-hiding
condition of Lemma~\ref{lem:htran-lhl} becomes
\[
  \ell_{r_0}\log_2(2\beta_{r_0}+1)
  \ge
  \log_2 q + 2\lambda/N
  \approx 20.2359.
\]
Table~\ref{tab:lhl-design-points} records representative design points.  The parameter \(\beta_{r_0}\) is the
coefficient bound of the long linear-side randomizer, and the membership degree is the degree of the vanishing
polynomial used to certify coefficient membership in \([ -\beta_{r_0},\beta_{r_0}]\).

\begin{table}[H]
\centering
\renewcommand{\arraystretch}{1.15}
\begin{tabular}{@{}cccl@{}}
\toprule
\(\beta_{r_0}\) & \(\ell_{r_0}\) & membership degree & \textbf{Comment} \\
\midrule
5  & 6 & 11 & balanced choice; moderate witness width and moderate degree \\
8  & 5 & 17 & saves one ring element in the witness; slightly higher degree \\
31 & 4 & 63 & entropy-rich but likely too expensive in proof degree \\
1  & 13 & 3 & ternary-style alphabet; degree is small but witness width is large \\
\bottomrule
\end{tabular}
\caption{Representative target-hiding design points for the long BB-tran randomizer.  The LHL condition constrains only
$r_0$; it does not change sampler tuning or the admissibility budget for $r_1$.}
\label{tab:lhl-design-points}
\end{table}

The table shows the main trade-off induced by the new target-hiding theorem.  Increasing \(\beta_{r_0}\) reduces the
required width \(\ell_{r_0}\), hence the number of witness coordinates and linear constraints, but it also increases the
degree of the range-check polynomial.  In the present regime, \((\beta_{r_0},\ell_{r_0})=(5,6)\) and
\((\beta_{r_0},\ell_{r_0})=(8,5)\) are the most natural candidates.

\subsection{SmallWood soundness terms}
SmallWood’s soundness is a combination of degree soundness (DECS), LVCS binding, and PIOP sampling errors
\cite{EPRINT:FenRiv25b}. For parameters \((N,s,\ell,\ell',\rho,\eta)\) and degree bounds \(d_{\mathrm{decs}}=s+\ell-1\)
and \(d_Q\), typical upper bounds take the form:
\[
\varepsilon_{\mathrm{decs}}\ \lesssim\ \binom{N}{d_{\mathrm{decs}}+2}\,q^{-\eta},
\qquad
\varepsilon_{\mathrm{sum}} = q^{-\rho},
\qquad
\varepsilon_{\mathrm{tail}} \approx \frac{\binom{d_Q}{\ell'}}{\binom{|S|}{\ell'}}.
\]
Our issuance statement includes the bilinear inverse check \((B_3-r_1)\Had Z=1\) together with the linear target identity given public \(T\), so \(d_Q\) must account for both the inverse relation and the range/membership checks.  The
showing statement additionally depends on the PRF effective degree (Appendix~\ref{app:prf}).

\subsection{SmallWood knobs (example)}
For the issuance and showing proofs discussed in Section~\ref{sec:parameters}, one example parameter tuple is:
\[
  (s,\ell,\ell',\rho,\theta,\eta)\;=\;(16,25,2,2,4,19),
\]
with masking degree \(d_Q\) derived from the instantiated constraint degrees. The pre-sign statement includes the bilinear inverse-witness constraint together with linear commitment, centering, and target equations, while the post-sign statement adds the PRF S-box composition on top of the same BB-tran inverse check.

\subsection{Extended benchmark grid (SmallWood)}
For additional context on how SmallWood knobs affect proof size and proving time, Table~\ref{tab:results-grid} reports
an example sweep over packing and soundness parameters (small-field variant).

\begin{table}[H]
\centering
\footnotesize
\setlength{\tabcolsep}{4pt}
\begin{tabular}{r r r r r r r r r}
\toprule
\textbf{ProofKB} & \textbf{Time (s)} & \textbf{Bits} & \textbf{NCols} & \(\ell\) & \(\ell'\) & \(\rho\) & \(\theta\) & \(\eta\)\\
\midrule
\ \ 9.32 & 0.266 & 133.44 & 6 & 22 & 2 & 2 & 4 & 17 \\
10.07 & 0.228 & 146.93 & 4 & 24 & 2 & 2 & 4 & 17 \\
11.52 & 0.250 & 133.22 & 4 & 28 & 2 & 2 & 4 & 17 \\
\midrule[1.2pt]
14.06 & 0.630 & 198.01 & 4 & 34 & 8 & 5 & 2 & 26 \\
14.17 & 0.654 & 192.26 & 6 & 34 & 8 & 5 & 2 & 26 \\
14.24 & 0.670 & 198.48 & 4 & 34 & 8 & 6 & 2 & 26 \\
\bottomrule
\end{tabular}
\caption{Example sweep over SmallWood knobs (proof size, prover time, and soundness bits). Sizes exclude public
parameters.}
\label{tab:results-grid}
\end{table}
